\section{Baby PIH}
\label{sec:baby}
\noindent In this section, we analyze the direct product construction to prove Baby PIH under $\wone \neq \fpt$.
\begin{theorem}[Main]\label{thm:baby}
    For any integer $r>1$, there is an integer $t>0$ such that for any 2CSP instance $\Pi=(X,\Sigma,\Phi)$ and its $t$-wise direct product instance $\Pi^{\odot t}=(X',\Sigma',\Phi')$:
    \begin{itemize}
        \item (Completeness) If $\Pi$ is satisfiable, then $\Pi^{\odot t}$ is satisfiable as well.
        \item (Soundness) If $\Pi$ is not satisfiable, then $\Pi^{\odot t}$ is not $r$-list satisfiable.
    \end{itemize}
\end{theorem}

Since $t$ is a constant depending solely on $r$, the number of variables in the new instance $|X'|=\binom{|X|}{t}$ depends only on $|X|$ rather than $|\Sigma|$, and the alphabet size of the new instance $|\Sigma'|\le |\Sigma|^t$ is polynomial in $|\Sigma|$. Thus for any constant $C$ and any computable function $f$, $f(|X'|)\cdot |\Sigma'|^{C}$ time is upper bounded by $g(|X|)\cdot |\Sigma|^{C\cdot t}$ time for some computable function $g$. Therefore, we have the following corollary from Theorem \ref{thm:baby}:
\begin{corollary}
    Assuming $\wone\neq\fpt$, Baby PIH is true.
\end{corollary}

Note that the completeness in~\Cref{thm:baby} follows trivially by assigning the restriction of the satisfying assignment on $X$ to each subset of variables. The main challenge is to show that when $\Pi^{\odot t}$ has a $r$-list satisfying assignment, $\Pi$ is satisfiable. To prove this, we will first work on a bipartite version of the direct product of 2CSP. 
\begin{definition}[Bipartite Direct Product Construction]
    Given a 2CSP instance $\Pi=(X,\Sigma,\Phi)$ and positive integers $a,b$, the $(a,b)$-bipartite direct product 2CSP instance $\Pi^{\odot (a,b)}=(X',\Sigma',\Phi')$ is constructed as follows.
    \begin{itemize}
        \item The variable set $X'$ consists of all $a$-sized subsets of $X$ on the left side, and all $b$-sized subsets of $X$ on the right side. With a little abuse of notation, we have $X' = \binom{X}{a} \cup \binom{X}{b}$.
        \item The domain of each variable $S \in X'$ is the set of all partial satisfying assignments for $S$ in $\Pi$. %The alphabet $\Sigma'$ is equal to $\binom{\Sigma}{a} \cup \binom{\Sigma}{b}$.
        \item For every $S \in \binom{X}{a}$ and $ T \in \binom{X}{b}$, we have a constraint in $\Phi'$ that checks whether $\sigma_S$ and $\sigma_T$ are consistent on the values for variables in $S \cap T$. 
    \end{itemize} 
\end{definition}

If for a 2CSP instance $\Pi$, $\Pi^{\odot t}$ is $r$-list satisfiable for $t=\max(a,b)$, by taking restrictions of its assignments on the smaller sets, it is easy to see $\Pi^{\odot(a,b)}$ is $r$-list satisfiable as well.
Thus, our goal is to show that if the bipartite instance $\Pi^{\odot(a,b)}$ is $r$-list satisfiable, then the original instance $\Pi$ is satisfiable.  

Our proof idea is borrowed from the Baby PCP theorem recently proved by Barto and Kozik \cite{BK22}. However, their theorem crucially relies on the fact that the alphabet $|\Sigma|$ is as small as a constant, which helps them to extract satisfying assignments for the smaller instance. The running time of their reduction is, therefore, $|X|^{O_{|\Sigma|}(1)}$, which is not affordable here since $|\Sigma|$ is as large as the input size. We resolve this issue by making use of the structural properties of the assignments in the direct product construction arising from the fact that they satisfy $r$-list consistency. 
If we fix some set $S$ on one side and consider its $r$ assignments, each set on the other side that intersects $S$ must have one of the $r$, which is a constant, assignments for their intersection part. We use this simple yet very useful observation when extracting the assignments in the inductive proof. 


In the following, we first prove Lemma \ref{lemma:2}, which is crucial to extract list satisfying assignments for the smaller subsets. Then in Lemma \ref{lemma:3}, we analyze a special case of the bipartite direct product construction when each variable on the right (bigger) side has only one assignment, and the consistency requirement is a slightly weaker one. In Lemma \ref{lem:final-baby}, we finish the analysis of the bipartite direct product construction, from which we get~\Cref{thm:baby} as a corollary.

\begin{lemma}\label{lemma:2}
    Let $k,r,q,a,b,b'$ be integers satisfying $r,q>0, a \ge k, b \ge r\cdot b'+a$. Consider the $(a,b)$-bipartite direct product 2CSP instance based on $\Pi=(X,\Sigma,\Phi)$. Let $u$ be an $r$-sized multi-assignment for $\binom{X}{a}$ and $v$ be a $q$-sized multi-assignment for $\binom{X}{b}$. Suppose for every $S \in \binom{X}{a}, T \in \binom{X}{b}$ with $T \supseteq S$, $v(T)_{|S} \cap u(S) \ne \emptyset$. 
    Then for every $A \in \binom{X}{k}$, there is an assignment $f_A$ for $A$ such that for every $T' \in \binom{X}{b'}$, there is some $T \in \binom{X}{b}$ satisfying $T \supseteq T' \cup A$ and $v(T)_{|A} \ni f_A$.
\end{lemma}
\begin{proof}
Suppose for the sake of contradiction that there is no such $f_A$ for some set $A \in \binom{X}{k}$. In other words, for any assignment $f$ on $A$, there exists a set $T'_{f} \in \binom{X}{b'}$, such that for every $T \in \binom{X}{b}$ satisfying $T \supseteq T'_{f} \cup A$, $v(T)_{|A} \not\ni f$.

Pick an arbitrary $S \in \binom{X}{a}$ with $S \supseteq A$. This can be done since $a \ge k$.  Consider any set $T \in \binom{X}{b}$ which contains $\cup_{f \in u(S)_{|A}} (T'_f \cup S)$. Such a $T$ exists since $b \ge r\cdot b'+a$. By the assumption about $A$, $v(T)_{|A}$ does not contain any value in $u(S)_{|A}$.  This contradicts the consistency guarantee in the hypothesis of the lemma, namely that for any $T \in \binom{X}{b}$ with $T \supseteq S$, $v(T)_{|S}$ must contain some value in $u(S)$. 
\end{proof}

See Figure \ref{fig:1} for an illustration of Lemma \ref{lemma:2}.

 \begin{figure}[h]
\centering
\includegraphics[height=9cm]{Baby_PIH/contents/fig1.png}
\caption{An illustration of Lemma \ref{lemma:2}.}
\label{fig:1}
\end{figure}

\begin{lemma}\label{lemma:3}
   For any integer $r > 0$, let $a=r$ and $b=(2r)^r$. Consider the $(a,b)$-bipartite direct product 2CSP instance based on $\Pi=(X,\Sigma,\Phi)$. Let $u$ be an $r$-sized multi-assignment for $\binom{X}{a}$ and $v$ be a 1-sized assignment for $\binom{X}{b}$. Suppose for every $S \in \binom{X}{a}$ and $T \in \binom{X}{b}$ with $S \subseteq T$, $v(T)_{|S} \in u(S)$. Then there is a global satisfying assignment $\sigma$ to the 2CSP instance $\Pi$.
\end{lemma}

\begin{proof}

We apply induction on $r$. When $r=1$, we have $a=1,b=2$, and both $u,v$ are 1-sized assignments. We claim that $u$ is the desired global satisfying assignment of $\Pi$. For each constraint on $(x,y) \in X^2$, $v(\{x,y\})=(x \mapsto u(x), y \mapsto u(y))$ by our consistency guarantee. By the construction of bipartite direct product 2CSP instance, the domain for each $T \in \binom{X}{b}$ consists only of partial satisfying assignments. Thus the fact that $(x \mapsto u(x),y \mapsto u(y))$ lies in the domain of $T=\{x,y\}$ implies $u$ satisfies the constraint on $(x,y)$.

When $r>1$, the idea is to extract consistent assignments for the $(a',b')$-bipartite direct product 2CSP instance and to decrease $r$, for some $a'\le a, b' \le b$. At a high level, if for some $x \in X$, every set in $\binom{X}{a}$ has at least two different values for it under $u$, we can keep only one of them and peel the other off to decrease $r$ by 1; otherwise we can prove the unique assignments are already satisfying.

In the following, define $k=1$, and $a'=a-1,b'=(2(r-1))^{r-1}$, i.e., $a',b'$ are parameters with respect to $r-1$. It's easy to see $k \le a$ and $r\cdot b'+a \le r\cdot (2r)^{r-1}+r \le (2r)^{r}=b$. According to Lemma \ref{lemma:2}, for every $A=\{x_i\}\in \binom{X}{k}$, there is an assignment $f_A$ for $A$ such that for every $T' \in \binom{X}{b'}$, there is some $T\in \binom{X}{b}$ satisfying $T \supseteq T' \cup A$ and $v(T)_{|A}\ni f_A$. Now $v(T)$ is of size-1, so we can simply write $v(T)_{|A}=f_A$.

Consider the following two cases:

\medskip\noindent
\textbf{Case 1:} For some $A=\{x_i\}$, Lemma \ref{lemma:2} holds for different assignments $p,q$. In other words, for every $T' \in \binom{X}{b'}$, there are $T_1,T_2 \in \binom{X}{b}$ satisfying $T_1,T_2 \supseteq T' \cup \{x_i\}$ and $v(T_1)_{|\{x_i\}}=p,v(T_2)_{|\{x_i\}}=q$. 
    
    Since for $r\ge 2$ we have $a=r\le (2(r-1))^{r-1}=b'$, for each $S \in \binom{X}{a}$ containing $x_i$, we can pick an arbitrary $T'\supseteq S$ and consider the corresponding sets $T_1,T_2$ above. By the consistency assumption between $(S,T_1)$ and between $(S,T_2)$ in Lemma \ref{lemma:3}, we can infer $u(S)_{|\{x_i\}}\ni p,q$.

    We construct new assignments $u',v'$ for the $(a',b')$-bipartite direct product 2CSP instance of $\Pi$, such that they still meet the consistency requirements in Lemma \ref{lemma:3}, and the size of $u'$ is at most $r-1$. For each $S' \in \binom{X}{a'}$, $u'(S')$ will be inherited from $u(S)$ for some $S \in \binom{X}{a}$ satisfying $S \supseteq S'\cup\{x_i\}$. Similarly for each $T' \in \binom{X}{b'}$, $v'(T')$ will be inherited from $v(T)$ for some $T \in \binom{X}{b}$ satisfying $T \supseteq T'\cup\{x_i\}$. 

    Suppose there is a arbitrary fixed order of all variables in $X$. We construct $S$ from $S'$ as follows, and output $u'(S')=\{\sigma \in u(S)|\sigma(x_i)=p\}_{|S'}$:
    \begin{itemize}
        \item If $x_i \not\in S'$, let $S=S'\cup \{x_i\}$.
        \item If $x_i \in S'$, $S$ is obtained by adding the lexicographical smallest element not in $S'$ to $S'$.
    \end{itemize}
    Note that we only keep the assignments in $u(S)$ whose restriction on $\{x_i\}$ is $p$, and discard those whose restriction on $\{x_i\}$ is $q$. Thus $u'$ is of size at most $r-1$ as desired.

    For each $T' \in \binom{X}{b'}$, we first construct $T''\in \binom{X}{b'}$ as follows:
    \begin{itemize}
        \item If $x_i \notin T'$, simply let $T''=T'$.
        \item If $x_i \in T'$, $T''$ is obtained by adding the lexicographical smallest element not in $T'$ to $T'$, and delete $x_i$.
    \end{itemize}
    After that, we find a $T \in \binom{X}{b}$ satisfying $T \supseteq T'' \cup \{x_i\}$ and $v(T)_{|\{x_i\}}=p$, and output $v'(T')=v(T)_{|T'}$.

    By our construction, $S' \subseteq T'$ implies $S \subseteq T$, so the new instance still meets the required consistency requirements, and the induction proceeds.

\medskip\noindent\textbf{Case 2:} Suppose for $A$ being each $\{x_i\}$, Lemma \ref{lemma:2} holds for a unique assignment $z_i$. 

    We claim that $\left(x_i \mapsto z_i\right)_{x_i \in X}$ is a global satisfying assignment. 
    Indeed, consider an arbitrary 
     constraint $\psi$ between variables $(x_i,x_j) \in X^2$. Applying Lemma~\ref{lemma:2} to the choice $A=\{x_i,x_j\}$, we know there is an assignment $f_A$ such that for every $T' \in \binom{X}{b'}$, there is some $T \in \binom{X}{b}$ satisfying $T \supseteq T' \cup A$ and $v(T)_{|A}=f_A$. Since $v(T)$ satisfies all constraints within $T$, this means that $f_A$ must satisfy the constraint $\psi$. By the uniqueness assumption of this case, we must have $f_A(x_i)=z_i$ and $f_A(x_j) = z_j$, which means that the assignment $(z_i,z_j)$ satisfies the constraint $\psi$ between $(x_i,x_j)$. 
     %\qedhere
\end{proof}

We will now finish the analysis of the bipartite direct product construction.
\begin{lemma}
\label{lem:final-baby}

For integers $r,q>0$, let $a=(2r)^{r+2q}, b=(2r)^{r+2q+2}$. Consider the $(a,b)$-bipartite direct product 2CSP instance based on $\Pi=(X,\Sigma,\Phi)$. Let $u$ be an $r$-sized multi-assignment for $\binom{X}{a}$ and $v$ be a $q$-sized multi-assignment for $\binom{X}{b}$. If $u,v$ list-satisfy all bipartite constraints, then there is a global satisfying assignment $\sigma$ to  the 2CSP instance $\Pi$.

\end{lemma}

\begin{proof}
We apply induction on $q$, the size of the right multi-assignment. When $q=1$, we have $a=(2r)^{r+2}\ge r$ and $b=(2r)^{r+4} \ge (2r)^r$. We can extract consistent satisfying assignments for the $(r,(2r)^r)$-bipartite direct product 2CSP instance, and invoke Lemma \ref{lemma:3} to prove there is a global satisfying assignment.

When $q>1$, we categorize discussions based on whether the left multi-assignment $u$ satisfies a certain property. Either we can still extract consistent assignments for the smaller $(a',b')$-bipartite direct product 2CSP instance while decreasing $q$ and leaving $r$ unchanged, or we can build multi-assignments that satisfy requirements in Lemma \ref{lemma:3}, and therefore directly invoke that lemma and stop the induction.

Define $k=(2r)^r$. The parameters $a',b'$ with respect to $q-1$ would be $a'=(2r)^{r+2q-2}$ and $b'=(2r)^{r+2q}$. In our case-analysis we will use the following inequalities, which we first prove here.

\begin{itemize}
    \item $b \ge r\cdot b'+a$.  Indeed 
    \[
    \begin{aligned}
        r\cdot b'+a & = r\cdot (2r)^{r+2q}+(2r)^{r+2q} \\
        & \le 2\cdot (2r)^{r+2q+1} \\
        & \le (2r)^{r+2q+2} =b.
    \end{aligned}
    \]
    \item $a \ge (r+1)\cdot (a'+k)$. Indeed
    \[
    \begin{aligned}
        (r+1)\cdot (a'+k) & = (r+1)\cdot ((2r)^{r+2q-2}+(2r)^r) \\
        & \le (2r)\cdot (2 \cdot (2r)^{r+2q-2}) \\
        & \le (2r)^{r+2q} = a.
    \end{aligned}
    \]
\end{itemize}

Now we consider the following two cases based on $u$. The criterion is whether $u$ satisfies a property which is reminiscent of the result of Lemma \ref{lemma:2}, except in Lemma \ref{lemma:2} the property is for $v$ while here we check it for $u$.

\begin{enumerate}
    \item Suppose there exists $A \in \binom{X}{k}$ such that for every assignment $f_A$ on $A$ and every $S' \in \binom{X}{a'}$, there exists $S \in \binom{X}{a}$ satisfying $S \supseteq S' \cup A$ and $u(S)_{|A} \not\ni f_A$.

    Given $a \ge k$ and $b \ge r\cdot b'+a$, we can apply Lemma \ref{lemma:2} to know that there is an assignment $f_A$ on $A$, such that for every $T' \in \binom{X}{b'}$, there is $T \in \binom{X}{b}$ satisfying $T \supseteq T'\cup A$ and $v(T)_{|A}\ni f_A$.

    We therefore build multi-assignments $u',v'$ for the smaller $(a',b')$-bipartite 2CSP instance as follows.

    For every $S' \in \binom{X}{a'}$, we pick the set $S$ guaranteed by the assumption of this case, and define $u'(S')=u(S)_{|S'}$. Note that $u'$ still has size at most $r$.

    For every $T' \in \binom{X}{b'}$, we pick the set $T$ guaranteed by Lemma \ref{lemma:2}, and define $v'(T')=\{\sigma \in v(T)| \sigma_{|A} \neq f_A\}_{|T'}$. By the assumption, there exists $S \in \binom{X}{a}$ satisfying $S \supseteq A$ and $u(S)_{|A}\not\ni f_A$. Thus to be consistent with $u(S)_{|A}$, we can conclude for every $T \in \binom{X}{b}$ containing $A$, $v(T)_{|A}$ should also contain some value other than $f_A$. Therefore, our constructed $v'$ is non-empty and has size at most $q-1$.

    It's also easy to see $u',v'$ still satisfy list consistency, since in $v'$ we only discard the assignments whose restriction on $A$ equals to $f_A$, which are not consistent with any $u'$.
    
    \item Suppose for every $A \in \binom{X}{k}$, there is an assignment $f_A$ for $A$ and a set $S' \in \binom{X}{a'}$, such that for every $S \in \binom{X}{a}$ satisfying $S \supseteq S' \cup A$, we have $u(S)_{|A}\ni f_A$.

    In this case we can construct $r$-sized multi-assignment $u'$ and $1$-sized assignment $v'$ for left and right part of the $(r,k=(2r)^r)$-bipartite direct product 2CSP of $\Pi$ respectively, that meets the requirements of Lemma \ref{lemma:3}. Furthermore, $u',v'$ are built purely based on $u$.

    For every $B \in \binom{X}{r}$, we define
    \[
    u'(B)=\bigcup_{A \in \binom{X}{k},A \supseteq B} (f_A)_{|B}
    \]
    For every $A \in \binom{X}{k}$, we define $v'(A)=f_A$.
    
    We first claim the size of $u'$ is at most $r$. Suppose it is not, there are $r+1$ sets $A_1,\ldots,A_{r+1} \in \binom{X}{k}$ with $f_{A_i}$ all different. Let $S'_1,\ldots,S'_{r+1}\in\binom{X}{a'}$ be the corresponding sets guaranteed in the assumption of this case. Consider a $S \in \binom{X}{a}$ which contains $\cup_{i=1}^{r+1}(S'_i \cup A_i)$. Such $S$ exists since $a \ge (r+1)\cdot (a'+k)$. Thus, $u(S)_{|A}$ contains $(r+1)$ different values, contradicting the fact that $u$ is of size $r$.

    It's easy to see $u',v'$ meets the requirement of Lemma \ref{lemma:3} by the definition of $u'$. Thus using Lemma \ref{lemma:3}, there is a satisfying assignment to the original instance $\Pi$. \qedhere
\end{enumerate}
\end{proof}

\noindent
Our main result, Theorem \ref{thm:baby}, now follows immediately from~\Cref{lem:final-baby}.

\begin{proof}[Proof of Theorem~\ref{thm:baby}]
    Given an integer $r$, we pick $t=(2r)^{r+2r+2}$, and prove that if a 2CSP instance $\Pi=(X,\Sigma,\Phi)$ is satisfiable, then so is $\Pi^{\odot t}$; if $\Pi$ is not satisfiable, then $\Pi^{\odot t}$ is not $r$-list satisfiable.

    The completeness case is easy: let $\sigma$ be a satisfying assignment for $\Pi$, then we can assign each set $S \in \binom{X}{t}$ the function that maps any $x \in S$ to $\sigma(x)$.

    For soundness case, suppose $\Pi^{\odot t}$ is $r$-list satisfiable by an assignment $\sigma$, we take $a=(2r)^{r+2q}, b=(2r)^{r+2q+2}$ and build multi-assignments $u,v$ for the left and right part of the $(a,b)$-bipartite direct product 2CSP instance $\Pi^{\odot(a,b)}$. For each set $S' \in \binom{X}{a}$, we pick an arbitrary $S \in \binom{X}{t}$ with $S \supseteq S'$, and define $u(S')=\sigma(S)_{|S'}$. Similarly for each $T'\in \binom{X}{b}$, we pick an arbitrary $T \in \binom{X}{t}$ with $T \supseteq T'$, and define $v(T')=\sigma(T)_{|T'}$. It's easy to see $u$ and $v$ are $r$-list consistent. Thus by Lemma \ref{lem:final-baby}, $\Pi$ is satisfiable.
\end{proof}